\chapter{Pensar en sistemas}
\opening{Comprender un avión es comprender relaciones entre componentes.}

En un capítulo técnico conviene separar narración y recuadros didácticos.

\begin{technicalbox}[Precaución editorial]
Esta obra es una guía educativa no oficial. No sustituye el manual de vuelo aprobado, las listas de verificación aplicables ni la instrucción aeronáutica certificada.
\end{technicalbox}

\section{Un ejemplo de ilustración}

\widefigure[0.8\linewidth]{assets/wing.pdf}{Esquema conceptual de demostración, no apto para instrucción operacional.}

La ilustración se guarda como archivo independiente y se reutiliza en distintos capítulos. La herramienta se puede adaptar a libros narrativos con ilustraciones ocasionales, manuales ilustrados y guías comerciales.

\sidecite[Una nota bibliográfica editorial]{Tufte, 2001}Una referencia lateral ayuda a ofrecer contexto sin convertir cada página en un artículo académico.

\section{Una idea final}

La plantilla distingue entre el contenido, que puede crecer cientos de páginas, y las decisiones de maquetación, que pueden cambiar antes de imprimir.
